\section{财产犯罪 - 基本原理、侵占罪}
\centering
\begin{forest}
  arrow to/.style n args=4{%
    delay={%
      tikz+={%
        \draw [every edge, line] () -- (!#1) node [above, midway] {#2};
      },
    },
  },
  for tree={
    grow'=east,
    edge={line},
    font=\small
  }
  [财产犯罪
    [非法占有目的
      [取得型犯罪
        [转移占有吗
          [转移，谁
            [犯罪人
              [暴力程度
                [抢劫]
                [抢夺]
                [盗窃]
              ]
            ]
            [被害人
              [诈骗]
              [敲诈勒索]
            ]
          ]
          [不转移
            [侵占]
          ]
        ]
      ]
      [故意毁坏财物]
    ]
  ]
\end{forest}
\begin{enumerate}
  \item 最重要，分论难度最大。体系思维，一棵树上结的果子。
  \item 非法占有目的（最最最最核心）\yjt 没有 \yjt 毁坏财物。有 \yjt 取得型犯罪。
        \begin{enumerate}
          \item 非法目的占有有时候题干会表述得很隐晦。\al{甲乙丙丁共谋诱骗戊参赌}。
          \item \bs{起意的阶段不同，结论也不同，这里会挖陷阱}。
        \end{enumerate}
  \item 谁实施了转移占有。
        \begin{enumerate}
          \item 犯罪人 \yjt 盗窃，位阶关系，转移手段\zd{暴力程度}不同。
                \begin{enumerate}
                  \item 盗窃。对物平和，对人无危险。
                  \item 抢夺。对物暴力，对人多多少少有危险。
                  \item 抢劫。对人暴力，压制反抗。
                \end{enumerate}
          \item 被害人 \yjt 诈骗、敲诈勒索。
        \end{enumerate}
  \item 保护法益。保护\bs{财产的秩序}，是民法的各种权的前提。
        \begin{enumerate}
          \item \zd{所有权}。
          \item \zd{占有事实}。包括非法占有（不允许黑吃黑）。
                \begin{enumerate}
                  \item \al{甲盗窃乙的摩托车放到后院，又被丙偷走}。
                  \item \al{甲盗回被乙偷走的摩托车}。偷回自己东西，是在行使所有权 \yjt 所以法条规定了“非法占有为目的”。
                  \item \al{甲提前偷回借出去的东西}。合法占有是所有权的基础，能够对抗所有权。没有非法所有目的，但有非法占有目的。
                \end{enumerate}
        \end{enumerate}
  \item 行为对象。财物，不能包括人的身体。
        \begin{enumerate}
          \item 例外。分离的人的身体器官。\al{如血液}。
          \item 尸体要具体问题具体分析。\al{大体老师、阴婚盗尸体}。
        \end{enumerate}
  \item \bs{非法占有目的}。\zd{非常重要，和后面贪污、受贿联系紧密}。
        \begin{enumerate}
          \item \bs{排除意思}/不具有返还意思。排除主人永久占有（\zd{不想还}）。
                \begin{enumerate}
                  \item \al{偷邻居家摩托车，就是为了去打个酱油}。
                  \item \al{晚上偷牛去犁地然后清晨还回去}。
                \end{enumerate}
          \item \bs{利用意思}。满足他的“需要”即可。
                \begin{enumerate}
                  \item 想用。\zd{毁掉不算}。\al{邻居敲木鱼太吵了，偷出来烧掉}。
                  \item \zd{不需要完全按照正常用途用}。\al{男偷女内衣、木质家具燃烧取暖、钢材当废品卖}。
                  \item 为自己为他人都可。\al{盗苹果手机送女朋友}。
                  \item 短暂利用也构成。\al{甲杀害后乙，盗走乙的车，开到河边沉车沉尸 \yjt 毁车行为属于不可罚的事后行为}。
                \end{enumerate}
        \end{enumerate}
  \item 侵占。\bs{他人所有自己占有 \yjt 自己所有}。
        \begin{enumerate}
          \item 对象。
                \begin{enumerate}
                  \item 代为保管物。\al{甲把钱给中间人乙去行贿官员，结果乙据为己有}。
                        \begin{enumerate}
                          \item 肯定说。\yi 符合构成要件。\er 不允许黑吃黑。
                          \item \zd{否定说}。钱性质改变 \yjt \zd{丧失了基于所有权的返还请求权，不构成侵占}。
                        \end{enumerate}
                  \item 埋藏物、遗失物。\yi 遗忘（临时/永久）。\er 场所如公园、公交车。
                  \item 列举不完，记行为公式而不是常见情形。
                \end{enumerate}
          \item 行使所有权。\zd{作为、不作为方式均可}。
                \begin{enumerate}
                  \item 特定物。\zd{据为己有、擅自处分}。
                  \item 种类物。现金。\zd{混同、拒不返还}。
                \end{enumerate}
        \end{enumerate}
\end{enumerate}
